%! TeX program = xelatex
\documentclass[../main.tex]{subfiles}

\begin{document}
\makelectureweek{8}{Week 8}

\section{L'H\^opital's Rule}
L'H\^opital's rule can be used to evaluate \emph{indeterminante forms}. We demonstrate its applications through examples and cautions. The precise theorem statement can be found in \href{https://eclass.yorku.ca/pluginfile.php/5312855/mod_resource/content/3/08_applications_2.pdf}{the weekly content pdf (week 8)} and page 310 of the textbook. 

\medskip

A limit \(\lim_{x \to a} \frac{f(x)}{g(x)}\) is an \emph{indeterminate form}
\begin{itemize}
  \item of type \hspace{2cm} if both \(\lim_{x \to a} f(x)\) and \(\lim_{x \to a} g(x)\) are 
    \bigskip
  \item of type \hspace{2cm} if both \(\lim_{x \to a} f(x)\) and \(\lim_{x \to a} g(x)\) are 
\end{itemize}
\medskip

\begin{example}
  Evaluate \(\lim_{x \to 1} \frac{\ln(x)}{x - 1}\).
\end{example}
\vfill

\clearpage
\begin{example}
  Evaluate \(\lim_{x \to \infty} \frac{e^{x}}{x^{2}}\).
\end{example}
\vfill


In general, L'H\^opital's rule can be used to evaluate limits (approaching a number or an infinity) of types
\[
  \frac{0}{0}, \quad \frac{\infty}{\infty}, \hspace{1in} 0 \cdot \infty, \hspace{1in} \infty - \infty, \hspace{1in} \quad 0^{0}, \quad \infty^{0}, \quad 1^{\infty}.
\]
\vspace{1in}

\bigskip

% \clearpage
% 0/0 or oo/oo are foundation. 
% convert all other types to either 0/0 or oo/oo.

\begin{example}
  % Convert \(\lim_{x \to 0^{+}} x \ln(x)\) to one of \(\frac{0}{0}\) orssheirf@yorku.ca \(\frac{\infty}{\infty}\). Then apply L'H\^opital's Rule.
  Evaluate \(\lim_{x \to 0^{+}} x \ln(x)\).
\end{example}
\vfill
\clearpage

\begin{example}
  % Convert \(\lim_{x \to \infty} (e^{x} - x)\) to one of \(\frac{0}{0}\) or \(\frac{\infty}{\infty}\). Then apply L'H\^opital's Rule.
  Evaluate \(\lim_{x \to \infty} (e^{x} - x)\).
\end{example}
\vfill

\begin{example}
  Evaluate \(\lim_{x \to 0^{+}} \sqrt[x]{x}\).
\end{example}
\vfill

\clearpage
\section{Curve Sketching}
Why do we learn curve sketching if a computer can do it? 
\[
  \href{https://www.geogebra.org/calculator/cpjqfsen}{y = 6x^{5} - 521 x^{4} + 2122 x^{3} - 3300 x^{2}.}
\]

\vfill

We will sketch graphs of functions using the folloing information if exists.
\begin{enumerate}[label={\Alph*}.]
  \item Domain. 
    \hfill
    {\footnotesize (from prerequisite knowledge)}

    \medskip

  \item Intercepts if exist and can be easily obtained.
    \hfill
    {\footnotesize (from prerequisite knowledge)}

    \medskip

  \item Symmetry (even/odd and periodic) if applicable.
    \hfill
    {\footnotesize (from prerequisite knowledge)}

    \medskip

  \item Asymptotes if exist.
    \hfill
    {\footnotesize (Sep. 7 and 21)}

    \medskip

  \item Intervals of increase or decrease.
    \hfill
    {\footnotesize (Oct. 31)}

    \medskip

  \item Local maximum or minimum values.
    \hfill
    {\footnotesize (Oct. 31)}

    \medskip

  \item Concavity and inflection points.
    \hfill
    {\footnotesize (Nov. 2)}

\end{enumerate}
\bigskip

\faExclamationTriangle{} 
If we are asked to sketch a curve on an assessment and such information is not given, then we must calculate them before sketching a curve. 

\bigskip

Remember to
\begin{itemize}
  \item draw and label the axes,
  \item draw asymptotes with dashed lines,
  \item label intervals with relevant information, and
  \item label points on the curve using notation \((x,y)\).
\end{itemize}

\clearpage
\begin{tikzpicture}
  \draw[gray, very thin] (0,0) grid[step=0.5] (9,7);
\end{tikzpicture}

\vfill

\begin{tikzpicture}
  \draw[gray, very thin] (0,0) grid[step=0.5] (9,7);
\end{tikzpicture}

\vfill
\begin{tikzpicture}
  \draw[gray, very thin] (0,0) grid[step=0.5] (9,7);
\end{tikzpicture}
\clearpage

\begin{minipage}{.45\textwidth}
  \centering
  \includegraphics[scale=.8]{../standalones/build/plot_symmetry_x_squared}
\end{minipage}
\hfill
\begin{minipage}{.45\textwidth}
  \centering
  \includegraphics[scale=.8]{../standalones/build/plot_symmetry_cos}
\end{minipage}
\vfill

\begin{minipage}{.45\textwidth}
  \centering
  \includegraphics[scale=.8]{../standalones/build/plot_symmetry_x_cubed}
\end{minipage}
\hfill
\begin{minipage}{.45\textwidth}
  \centering
  \includegraphics[scale=.8]{../standalones/build/plot_symmetry_sin}
\end{minipage}
\vfill

\begin{minipage}{.45\textwidth}
  \centering
  \includegraphics[scale=.8]{../standalones/build/plot_periodic_sin}
\end{minipage}
\hfill
\begin{minipage}{.45\textwidth}
  \centering
  \includegraphics[scale=.8]{../standalones/build/plot_periodic_tan}
\end{minipage}
\vfill

\begin{minipage}{.45\textwidth}
  \centering
  \includegraphics[scale=.8]{../standalones/build/plot_no_symmetry_poly}
\end{minipage}
\hfill
\begin{minipage}{.45\textwidth}
  \centering
  \includegraphics[scale=.8]{../standalones/build/plot_no_symmetry_wavy}
\end{minipage}
\clearpage

\begin{tikzpicture}
  \draw[gray, very thin] (0,0) grid[step=0.5] (9,7);
\end{tikzpicture}

\vfill

\begin{tikzpicture}
  \draw[gray, very thin] (0,0) grid[step=0.5] (9,7);
\end{tikzpicture}

\vfill
\begin{tikzpicture}
  \draw[gray, very thin] (0,0) grid[step=0.5] (9,7);
\end{tikzpicture}
\clearpage

\phantom{top}
\vfill
\centerline{
  \begin{tikzpicture}
    \draw[gray, very thin] (0,0) grid[step=0.5] (18,10);
  \end{tikzpicture}
}
\clearpage

\section{Optimization Problems}

\begin{example}
  Consider a ball placed at height \(y = 0\). The ball is launched with initial velocity \(v_{0}\) and initial angle \(\theta\). 

  Ignore air resistance and assume gravity is the only force acting on the ball during flight. Then the horizontal range \(d\) is a function of the initial angle
  \[
    d(\theta) = \frac{v_{0}^{2}}{g} \sin(2 \theta).
  \]

  Find the initial angle \(\theta\) that maximizes the horizontal range \(d\).
\end{example}
\clearpage

\begin{example}
  Suppose we want to build a rectangular fence that encloses an area of \(100\) square feet. What dimension (length and width) minimizes the cost of the fence?
\end{example}
\clearpage

\phantom{in case we need more space.}
\clearpage

\begin{example}
  A haunted house is trying to maximize revenue for next year. This year, if the ticket price was \textdollar{20}, then the average attendance was \(500\) visitors per evening. If the ticket price was \textdollar{15}, then the average attendance was \(800\) visitors per evening. 

  Find the ticket price that maximizes revenue.
\end{example}
\clearpage

\phantom{in case we need more space.}
\clearpage
\end{document}
